\section{The actual tuple completion, presentation and collision cut}
\label{sec:cubic-tuple-local-presentation}

The original local presentation theorem is stated for a residual
character and its powers.  Here the residual datum is the actual
two-root cubic tuple.  We prove the completion and obstruction
interfaces for that datum, and then give the local collision estimate
with the joint residual inertia kernel.  The resulting component
bound does not by itself prove the parity presentation, the source
component-chaining assertion, or propagation along a prescribed curve.

Retain the actual host, groups, tuple and rings of AC-G94--AC-G99:
\[
 \begin{gathered}
 p=269,\quad \Gamma=G_F/N,\quad
 G=\ker(\eta_K:\Gamma\longrightarrow\{\pm1\}),\quad
 d=\sum_{v\mid p}d_v=[F:\mathbb Q],\\
 (\theta_0,\theta_1,\theta_2)
       =(\chi^j\gamma,1,\chi^j\gamma^{-1}).
 \end{gathered}
 \tag{G101.1}
\]
Here \(j\in\{1,2\}\), \(c\in\Gamma\) is complex conjugation,
\(g^\star=cg^{-1}c\), and
\(\gamma=\bar\psi\bar\epsilon^e\).
The actual local hypothesis is
\(\operatorname{comp}_p(\psi)^2=1\); no global assertion
\(\psi^2=1\) is made.  At the coefficient-prime places,
\(\bar\psi^2\) is unramified.  The initial completions of \(F\)
are unramified over \(\mathbb Q_p\), and the subsequent local
extensions have \(p\)-power degree.  Thus the image of
\(\bar\epsilon\) on inertia still has order \(268\).
The tame cubic character \(\chi\) still has order three there.
These are the actual local facts used below.

The complete ring \(R\), its closed weight-zero ring \(B\), and
\(A=B[r_{ij}(g)]\) are those of AC-G95.
Write \(x_0=\mathfrak m_BA\), and retain its notation
\[
 T_-=(r_{ij}(g):i<j),\qquad T_+=(r_{ij}(g):i>j).
 \tag{G101.2}
\]
Thus \(T_-=0\) is the lower-triangular condition.
At the normalization element \(g_0\),
\[
 r(g_0)-r(g_0^\star)=\operatorname{diag}(a,0,-a),
 \qquad \bar a=a_0\ne0.
 \tag{G101.3}
\]
Every local group has a separating element on which the three
residual characters have pairwise distinct values, as proved
in AC-G94.

Let \(x\) be a closed point of \(\operatorname{Spec}A\) over
\(\mathfrak m_B\).  Choose a finite unramified coefficient
extension \(\mathcal O'/\mathcal O\), a point \(y\) of
\(A'=\mathcal O'\otimes_{\mathcal O}A\) over \(x\), and an
identification of its residue field with \(k'=\mathcal O'/\varpi\).
Enlarging the unramified extension permits this choice.
Set
\[
                  C_x=\widehat{A'_y},\qquad
                  \sigma=r\bmod y .
 \tag{G101.4}
\]
The completion has the displayed coefficient map; all representations
and local characters below are in their specified bases.

\begin{lemma}[The actual tuple local-completion interface]
\label{lem:G101completion}
The ring \(C_x\) represents the continuous normalized polarized
ordinary lifts of \(\sigma\), with the same ordered local character
labels.  This assertion is valid for every complete Noetherian local
\(\mathcal O'\)-algebra with residue field \(k'\), including rings of
characteristic \(p\).
\end{lemma}

\begin{proof}
First \(\sigma\) is upper or lower triangular.  Indeed every
positive-weight entry and every negative-weight entry have a product
of positive powers of weight zero.  If both entries were nonzero
in the field \(k'\), that product would be nonzero.
But the corresponding universal weight-zero monomial belongs to
\(\mathfrak m_B\), and therefore vanishes at \(x\).
So all off-diagonal entries with one of the two signs vanish.
All diagonal entries and local characters at \(x\) are the fixed
residual values, by AC-G95.

We give the construction for a lower-triangular \(\sigma\).
The upper-triangular case uses the opposite diagonal rescaling.
Let \(D\) be a finite Artin local \(\mathcal O'\)-algebra with
residue field \(k'\), and let \(\rho\) be such a lift to \(D\).
For the moment write \(t\) for a formal variable.
The entries of
\[
 \rho^t(g)=
   \operatorname{diag}(t^{-1},1,t)\,
       \rho(g)\,
   \operatorname{diag}(t,1,t^{-1})
 \tag{G101.5}
\]
belong to \(D[t,t^{-1}]\).  Lower entries acquire positive powers
of \(t\); upper entries have coefficients in \(\mathfrak m_D\)
and acquire powers \(t^{-1}\) or \(t^{-2}\).
The diagonal and local-character residues are in \(k\subset k'\).
This rescaling preserves the star equation, the diagonal
normalization at \(g_0\), and the ordered ordinary product identity.

Here is a complete local test ring in which (G101.5) is a lift of
the original diagonal tuple over \(k\).
Inside \(D((t))\), put
\[
 L=D[[t]][\,\mathfrak m_Dt^{-2}\,],\qquad
 \mathcal N=\ker\bigl(L\longrightarrow k'[[t]]\bigr),\qquad
 L_0=\{z\in L:z\bmod(\mathcal N,t)\in k\}.
 \tag{G101.6}
\]
If \(\mathfrak m_D^\ell=0\), any product of \(\ell\)
generators with coefficients in \(\mathfrak m_D\) is zero.
It follows that \(L\) is a finite \(D[[t]]\)-module:
it is generated by the finitely many products of fewer than
\(\ell\) generators of \(\mathfrak m_Dt^{-2}\).
Its Laurent series have negative exponents bounded below.
Coefficient reduction gives the indicated surjection to \(k'[[t]]\).
Every coefficient of an element of \(\mathcal N\) is in
\(\mathfrak m_D\), so \(\mathcal N^\ell=0\).
In general \(\mathcal N\) is larger than
\(\mathfrak m_DL\); no equality with that smaller ideal is used.

The finite ring \(D\) is a finite \(\mathcal O\)-module.
Consequently \(L\) is finite over \(\mathcal O[[t]]\).
The subset \(L_0\) is an \(\mathcal O[[t]]\)-submodule, hence
is finite over this Noetherian ring as well, and it is a subring.
It is local with residue field \(k\):
an element with nonzero residue in \(k\) is invertible in \(L\),
and its inverse again has residue in \(k\).
Its maximal ideal is the preimage of \((\mathcal N,t)\).
Since \(L_0\) is finite over the complete local ring
\(\mathcal O[[t]]\), and that maximal ideal is the radical of
\((\varpi,t)L_0\), it is complete for its maximal-ideal topology.
Thus \(L_0\) is an allowed test ring of AC-G94.

The data (G101.5), with the original local characters, belong to
\(L_0\) and reduce to the diagonal tuple \(\boldsymbol\theta\).
Their inverse matrices also belong to \(L_0\), since they are
the same rescaling applied to \(\rho(g^{-1})\).
They are continuous: the original \(D\)-valued data have finite
image and open fibers, and rescaling preserves those fibers.
The universal property gives \(R\to L_0\).
Every zero-weight monomial and every local character maps to a
constant in \(D\).  The constants in \(L_0\) form a finite, and
therefore closed, subset of the Hausdorff ring \(L_0\).
Taking the defining closure of the weight-zero generators shows
that \(B\) maps to constants.
Since \(A\) is generated over \(B\) by actual entries, its image
is in \(D[t,t^{-1}]\).  Evaluation at \(t=1\) is legitimate
on this image and gives a homomorphism
\[
                         A\longrightarrow D
 \tag{G101.7}
\]
with the prescribed representation and characters.
The residue of (G101.7) is the given point \(x\).
It sends \(\mathfrak m_B\) into \(\mathfrak m_D\).

This map is unique among such maps.
Its restriction to \(B\) is continuous, since it is local.
The dense zero-weight generators and local characters determine
that restriction; the matrix entries then determine its values
on \(A\).  Combining (G101.7) with the given coefficient map
\(\mathcal O'\to D\), localizing at \(y\), and completing gives
the unique map \(C_x\to D\).

For a general complete Noetherian local target \(D\), apply the
finite-Artin result to every \(D/\mathfrak m_D^n\).
Uniqueness makes the maps compatible, and their inverse limit
is the required map \(C_x\to D\).  Conversely the universal data
on \(A\) are \(x_0\)-adically continuous by AC-G95, hence
continuous after passage to this completion and to any such
target.  All their polynomial identities survive.
This proves both directions of the claimed universal property.
\end{proof}

\paragraph{The local character ideals.}
Use the joint residual kernel, exactly as in AC-G102:
\[
 \begin{split}
 K_v&=I_v\cap\bigcap_{i=0}^2
                   \ker(\theta_i\circ\operatorname{loc}_v),\\
 N_v&=I_v\cap\ker\epsilon,\\
 \lambda_\theta&=(\varpi)B+
       (\,\kappa_{v,i}(g)-1:g\in K_v,\ v,i\,)B,\\
 \mathfrak c_{v,ij}&=
       (\,\kappa_{v,i}(g)-\kappa_{v,j}(g):g\in N_v\,)C_x
                         \quad(i\ne j).
 \end{split}
 \tag{G101.8}
\]
The \(N_v\) here uses the full cyclotomic character.
Let \([\theta_i]\) denote the multiplicative Teichmuller lift of
the finite residual character to \(\mathcal O^\times\).  Define also
\[
 \widehat\kappa_{v,i}(g)
       =\kappa_{v,i}(g)[\theta_i(\operatorname{loc}_v g)]^{-1},
 \qquad
 \widehat{\mathfrak c}_{v,ij}
       =(\,\widehat\kappa_{v,i}(g)-\widehat\kappa_{v,j}(g):
                                                    g\in N_v\,)C_x .
 \tag{G101.8a}
\]
These are collisions after removing the residual characters.
The Teichmuller lifts exist uniquely by simple-root Hensel lifting
of \(X^{\#k^\times}-1\), and form multiplicative characters.
On \(N_v\), the residual endpoint characters agree:
\(\theta_0/\theta_2=\gamma^2=1\).
Thus \(\widehat{\mathfrak c}_{v,02}=\mathfrak c_{v,02}\).
For an adjacent pair the two notions differ.
For an ideal \(I\subset B\), the required finite weight-fiber
condition is
\[
                \sqrt{I+\lambda_\theta}=\mathfrak m_B.
 \tag{G101.9}
\]
For the actual native support ideal
\(I=\Phi(J_{\rm aut})B\), AC-G102-finlam proves (G101.9)
under its exact retained native support and \texttt{ModU} inputs.
No additional arithmetic FinLam hypothesis is being silently added.

\begin{theorem}[AC-G101-local: the tuple presentation and collision bounds]
\label{AC-G101-local}
For every completion \(C_x\) in (G101.4), there are integers
\(N,r\ge0\) and a continuous surjection
\[
 \mathcal O'[[X_1,\ldots,X_N]]\longrightarrow C_x
 \tag{G101.10}
\]
whose kernel has \(r\) generators, with
\[
                         N-r\ge3d+1.
 \tag{G101.11}
\]
Every irreducible component of \(\operatorname{Spec}C_x\)
therefore has dimension at least \(3d+2\).
By padding the relation list with zero relations, (G101.10)
has the literal numerical form of the source's \texttt{Loc.Pres}.

Both adjacent collision ideals are the unit ideal:
\[
                  \mathfrak c_{v,01}
                    =\mathfrak c_{v,12}=C_x .
 \tag{G101.12}
\]
If \(I\subset B\) satisfies (G101.9), and a prime
\(\mathfrak p\subset C_x\) is off both triangular loci and
contains \(IC_x+\widehat{\mathfrak c}_{v,ij}\), for any
pair \(i\ne j\), then
\[
 \begin{split}
 \dim C_x/\mathfrak p+d_v&\le3d+3,\\
 \varpi\in\mathfrak p
       \quad\Longrightarrow\quad
 \dim C_x/\mathfrak p+d_v&\le3d+2.
 \end{split}
 \tag{G101.13}
\]

In particular suppose \(T_+\not\subset x\), so that the
actual AC-G98 triangular bound applies, and put
\[
                   D_{\rm tri}=2d+h+a+z+1 .
 \tag{G101.14}
\]
Then
\[
 \dim C_x/(IC_x+(\varpi)+\widehat{\mathfrak c}_{v,ij})
       \le\max\{D_{\rm tri},\,3d+2-d_v\}.
 \tag{G101.15}
\]
If \(h+a+z+1<d\) and \(d_v>2\), this dimension is less
than \(3d\).  Thus the local finite-image-ratio locus in
the characteristic-\(p\) native-support intersection has the
strict dimension cut needed at this step, for all three pairs.
In a characteristic-\(p\) prime quotient, vanishing of
\(\widehat{\mathfrak c}_{v,ij}\) is equivalent to the ratio
\(\kappa_{v,i}/\kappa_{v,j}\) having finite image on \(N_v\).
Outside this closed locus the two characters remain distinct
on every open subgroup of \(N_v\), and hence on every open
subgroup of the full local group.
The adjacent unit-ideal assertion (G101.12) alone would not
imply this stronger separation.
This does not assert a parity presentation or the remaining
component-propagation conclusions.
\end{theorem}

\begin{proof}
\emph{1. The actual adjoint cone and its Euler characteristic.}
Extend \(\sigma\) to the polarized semidirect group on
\(\Gamma=G\rtimes\langle c\rangle\), with fixed second coordinate
\(\eta_K\) and \(c\) represented by \((1,c)\).
The involution on matrices is
\(\Theta(M)=JM^{-t}J\).
Let \(W=M_3(k')\), with conjugation action on \(G\) and
action \(X\mapsto-JX^tJ\) at \(c\).
At each coefficient-prime place let \(\mathfrak b_v\subset W\)
be the stabilizer of the ordered flag of \(\sigma\).
It is a conjugate of the upper-triangular Lie algebra and has
dimension six.  Let \(H^i_{\mathfrak b}\) denote the actual
continuous cone cohomology
\[
 \operatorname{Cone}\left(
 C^\bullet(\Gamma,W)\longrightarrow
 \bigoplus_{v\mid p}C^\bullet(G_{F_v},W/\mathfrak b_v)
                    \right)[-1],
 \qquad h_i=\dim_{k'}H^i_{\mathfrak b}.
 \tag{G101.16}
\]

The local separating element shows that
\[
                  H^0(G_{F_v},W/\mathfrak b_v)=0.
 \tag{G101.17}
\]
In its eigenbasis the quotient has the strictly lower
matrix positions, whose eigenvalues are the nontrivial ratios
of the three distinct eigenvalues.  None is one.
This also proves the uniqueness of the ordered invariant flag
over an Artin lift, by the labeled Hensel projectors.

We check the dual-vanishing hypothesis of AC-G96-EC directly,
instead of importing the original one-character \texttt{pr\_dual}.
The residual \(\sigma|_G\) is triangular by the completion lemma.
The composition factors of its adjoint are
\(\theta_i\theta_j^{-1}\).
Let \(c_p\) denote the actual action on \(\mu_p\), which is
\texttt{chiO} in the duality source.  The notation \(\epsilon\)
elsewhere in this package may use the inverse geometric convention;
we therefore check both \(\bar\epsilon\) and
\(\bar\epsilon^{-1}\), without identifying them silently with \(c_p\).
The four adjacent factors have nontrivial cubic inertia,
whereas both cyclotomic conventions are trivial on tame pro-three inertia.
The three diagonal factors are trivial, whereas either
\(\bar\epsilon^{\pm1}\) has order \(268\) on inertia.
The two endpoint factors are \(\gamma^{\pm2}\); on inertia
they are \(\bar\epsilon^{\pm2e}\), since
\(\operatorname{comp}_p(\psi)^2=1\).
Neither can equal either \(\bar\epsilon\) or \(\bar\epsilon^{-1}\),
because each of \(\pm2e-1\) and \(\pm2e+1\) is odd and cannot
be divisible by \(268\).  Consequently
\[
                 \operatorname{Hom}_\Gamma
                         (W,k'(c_p))=0.
 \tag{G101.18}
\]
Here equality of global factors would imply equality on this
actual full local group; the local calculation is sufficient in
both arithmetic and inverse geometric cyclotomic conventions.

At any real place the action on \(W\) has the form
\(X\mapsto-SX^tS^{-1}\), where \(S\) is an invertible
symmetric matrix: the corresponding semidirect element has
order two.  Its invariant space is identified by \(X\mapsto XS\)
with the skew-symmetric matrices, of dimension three.
The actual Euler formula AC-G96-EC, with the finite-field
extension \(k'\), gives
\[
 h_2+h_0+6d=h_1+3d,\qquad
                   h_1=h_2+h_0+3d.
 \tag{G101.19}
\]
Its sheaf/cohomology and duality inputs are exactly those already
retained in AC-G96; they are not reproved or strengthened here.
In particular (G101.18) also gives vanishing in degrees
above two, so there is no uncounted degree-three term.

\emph{2. Tangents and the one remaining normalization direction.}
Let \(T\) be the tangent space of the functor in the completion
lemma, so
\[
       N=\dim_{k'}T
          =\dim_{k'}\mathfrak m_{C_x}/
                      (\mathfrak m_{C_x}^2+\varpi C_x).
 \tag{G101.20}
\]
A tangent vector, extended to the polarized semidirect group,
gives a continuous adjoint one-cocycle; variation of its local
flags gives the local cone coordinates.  Modulo global strict
conjugation its image is a class in \(H^1_{\mathfrak b}\).
This map is onto.  We spell out the normalization point.
An arbitrary cone cocycle constructs a representation over
\(k'[\delta]/\delta^2\), with its local invariant flags.
If its value at \(c\) is \((1+X,c)\), then
\(JX^tJ=X\), and conjugation by \(1-X/2\) changes that value
to \((1,c)\).  Next its normalization matrix differs from
\(\operatorname{diag}(a_0,0,-a_0)\) by a first-order
\(J\)-skew matrix.  Conjugating by \(1+Y\), with
\(JY^tJ=-Y\), kills its off-diagonal entries:
the off-diagonal commutator entries are
\((d_j-d_i)Y_{ij}\), for
\((d_0,d_1,d_2)=(a_0,0,-a_0)\), and the differences
are units.  The \(J\)-skew identities are preserved by this
solution.  Local flags are transported by the same conjugations.
The completion lemma then gives a tangent vector of \(C_x\).
This proves surjectivity up to the global coboundary, which is
exactly the asserted map to \(H^1_{\mathfrak b}\).

For completeness, (G101.17) is what prevents loss of a local
cone class when one forgets its chosen flag primitive.
If a cocycle is a global coboundary \(gX-X\), its local
primitive \(Y_v\) satisfies
\(g(X-Y_v)-(X-Y_v)\in\mathfrak b_v\).
Equation (G101.17) gives \(X-Y_v\in\mathfrak b_v\);
the entire cone class is consequently a boundary.

Put \(X_0=\operatorname{diag}(1,0,-1)\ne0\).
Conjugation by \(\operatorname{diag}(1+\delta,1,(1+\delta)^{-1})\)
preserves every normalization equation.  It gives a tangent
whose cocycle is, up to the harmless overall sign,
the coboundary \(gX_0-X_0\).
If \(X_0\) is not globally invariant, this is a nonzero tangent
in the kernel of the map to \(H^1_{\mathfrak b}\), so
\(N\ge h_1+1\).
If \(X_0\) is globally invariant, (G101.17) forces
\(X_0\in\mathfrak b_v\) for every \(v\), and therefore
\(h_0\ge1\), while surjectivity still gives \(N\ge h_1\).
In both cases (G101.19) yields
\[
                            N\ge h_2+3d+1.
 \tag{G101.21}
\]
This is the actual tuple version of the complete linear-algebra
argument in \texttt{PrCount.lean}.

\emph{3. A complete small-extension obstruction space.}
Let \(E\to D\) be a surjection of complete Noetherian local
\(\mathcal O'\)-algebras with residue field \(k'\), whose kernel
\(K\) is annihilated by \(\mathfrak m_E\).
Lift continuously the matrices of a \(D\)-valued representation
and the upper-triangular local matrices in bases for its local flags.
Such continuous set lifts exist also when \(D\) is not Artinian.
Indeed \(K\) is finite and
\(\bigcap_n\mathfrak m_E^n=0\); for large \(n\),
\(K\cap\mathfrak m_E^n=0\).
The map \(\mathfrak m_E^n\to\mathfrak m_D^n\) is then
a continuous bijection of compact Hausdorff additive groups,
and hence a homeomorphism.  Lift finitely many coset
representatives and use this section on each coset.
Invertible matrices and upper-triangular invertible matrices
lift because the kernel is nilpotent.

Write \(\widetilde\rho\) for a set lift into the semidirect
group with its fixed second coordinate, and
\(\widetilde\ell_v\) for the conjugate of a local
upper-triangular set lift by a lifted flag basis.
The two lifts agree modulo \(K\) on the local group.
Under the additive identification
\(1+M_3(K)\simeq W\otimes K\), define
\[
 \begin{split}
 c(s,t)&=\widetilde\rho(s)\widetilde\rho(t)
                       \widetilde\rho(st)^{-1}-1,\\
 a_v(s)&=\widetilde\rho(\operatorname{loc}_v s)
                       \widetilde\ell_v(s)^{-1}-1 .
 \end{split}
 \tag{G101.22a}
\]
The products in the first line are semidirect products; their
second coordinate is one.  The global multiplication defect
is thus a continuous two-cochain with values in \(W\otimes K\).
The difference between its restriction and the multiplication
defect of the chosen upper-triangular local lift is the
coboundary of a local one-cochain modulo \(\mathfrak b_v\).
In the usual additive notation, the pair has
\[
 dc=0,\qquad
       \operatorname{res}_v(c)-da_v
                    \in C^2(G_{F_v},\mathfrak b_v\otimes K).
 \tag{G101.22}
\]
These identities follow by expanding three factors in the two
associative orders; every product of two kernel entries is zero.
Changing set lifts adds a cone boundary.
Thus the defect determines
\[
                         o(\rho,E)\in H^2_{\mathfrak b}\otimes K,
 \tag{G101.23}
\]
functorially and linearly in the kernel.
If this class is zero, subtract the global one-cochain
in a cone primitive from the matrix lifts, and use its local
zero-cochains to conjugate the lifted local flags.
Equation (G101.22) then gives a global homomorphism preserving
those local flags.  Conversely any such lift makes the class zero.
The diagonal entries of each corrected local triangular
representation give the required local characters.

Finally normalize at \(c\) and \(g_0\) by the two strict
conjugations of Step 2.  The same formulas work over this
small extension, because \(K^2=0\) and
\(\mathfrak m_EK=0\); all commutator coefficients reduce to
the same unit differences \(a_0,2a_0\).
They do not alter the prescribed reduction.
Consequently (G101.23) is a complete obstruction space for
the actual normalized functor, not just for an unframed
or unnormalized auxiliary problem.

\emph{4. Relations and component dimension.}
Choose the minimal power-series surjection (G101.10), with
\(N\) as in (G101.20), and write \(S=\mathcal O'[[X_1,\ldots,X_N]]\),
\(\mathfrak a=\ker(S\to C_x)\).
Such a surjection follows by lifting a relative cotangent basis
and successive approximation in the complete topology.
Minimality says
\(\mathfrak a\subset\mathfrak m_S^2+\varpi S\).
The finite vector space
\[
                         K=\mathfrak a/\mathfrak m_S\mathfrak a
 \tag{G101.24}
\]
is the kernel of the complete small extension
\(S/\mathfrak m_S\mathfrak a\to C_x\).
Its obstruction gives a linear map
\(K^\vee\to H^2_{\mathfrak b}\).
This map is injective, as follows.
If a functional \(u\) has zero obstruction, push the extension
out by \(u\).  By Step 3 and the completion lemma, the resulting
extension \(E_u\to C_x\) has a section as a local
\(\mathcal O'\)-algebra map: the lifted universal representation
gives the section, and universality makes its composite the identity.
The natural map \(S\to E_u\) and the map through that section
differ by an \(\mathcal O'\)-linear derivation with values in
the residue-field kernel.
Such a derivation kills \(\mathfrak m_S^2+\varpi S\).
On \(\mathfrak a\) its value is \(u\), so \(u=0\).
It follows that \(\dim K\le h_2\).
Nakayama identifies \(\dim K\) with the minimal number \(r\)
of generators of \(\mathfrak a\).
Together with (G101.21), this proves (G101.11).

The ring \(S\) is regular local of dimension \(N+1\).
The generalized principal ideal theorem bounds the height of
every minimal prime over \(\mathfrak a\) by \(r\).
The dimension formula in this regular local ring then gives
dimension at least \(N+1-r\ge3d+2\) for every component.
If desired set \(r'=N-(3d+1)\ge r\), and append \(r'-r\)
zero relations to get the exact numerical
\texttt{Loc.Pres} equality.

\emph{5. The adjacent collisions are empty.}
Tame pro-three inertia contains an element \(s\) on which
\(\chi(s)\) has order three.
It belongs to \(N_v\), since the image of the full cyclotomic
character lies in
\(\mathbb Z_{269}^{\times}\), whose prime-to-\(269\) torsion
has order \(268\), and whose pro-\(269\) part has no
pro-three quotient.
The rational local character has order at most two on
this pro-three subgroup, so \(\gamma(s)=1\).
The residual local values are
\((\chi^j(s),1,\chi^j(s))\).
Both adjacent differences are therefore units in \(B\), and
remain units in \(C_x\).  This proves (G101.12).

\emph{6. Counting every residually normalized collision cut in \(B\).}
Let \(\mathfrak p\) be as in (G101.13) and
\(\mathfrak q=\mathfrak p\cap B\).
Use the \(d_w\) generators \(u_{w,a}\) of AC-G99:
they kill every residual tuple character, all except
\(u_{w,0}\) lie in \(N_w\), and every finite abelian
\(p\)-character which kills all these generators is unramified.

In \(B/\mathfrak q\), impose
\(\kappa_{w,i}(u_{w,a})-1\) for every triple \((w,i,a)\)
except the block with local place \(v\) and the second member
\(j\) of the colliding pair; in that omitted block impose
only \(\kappa_{v,j}(u_{v,0})-1\).
There are \(3d-d_v+1\) displayed equations.
Also impose \(\varpi\) if it is not already in
\(\mathfrak q\).  The normalized collision recovers all
remaining omitted equations, because \(u_{v,a}\in N_v\)
for \(a\ne0\), and every residual tuple character is one
on these generators.  Thus the normalized and unnormalized
character values at these generators coincide exactly.

We prove that the radical of this ideal is the maximal ideal
of \(B/\mathfrak q\).
Let \(\mathfrak q'\supset\mathfrak q\) be a prime containing
all the imposed equations.
Its quotient has characteristic \(p\).
For one local character \(\kappa\), let \(m=\#k^\times\).
On each finite quotient by a power of the maximal ideal,
\(\kappa^m\) is a continuous finite abelian \(p\)-character:
its residue is one and the principal-unit group of a finite
local characteristic-\(p\) ring is a finite \(p\)-group.
It kills all the \(u_{w,a}\).
AC-G99 therefore makes it trivial on inertia.
Passing through all powers gives
\(\kappa(g)^m=1\) on inertia.
If \(g\in K_w\), then \(\kappa(g)\) has residue one;
the factor
\(1+\kappa(g)+\cdots+\kappa(g)^{m-1}\)
has the nonzero residue \(m\) and is a unit.
Thus \(\kappa(g)=1\).
This proves \(\lambda_\theta\subset\mathfrak q'\).
Since \(I\subset\mathfrak q\), (G101.9) now forces
\(\mathfrak q'=\mathfrak m_B\).

All the cut equations lie in the maximal ideal.
The parameter-ideal form of the principal ideal theorem yields
\[
 \begin{split}
  \dim B/\mathfrak q&\le3d-d_v+2,\\
  \varpi\in\mathfrak q
       \quad\Longrightarrow\quad
  \dim B/\mathfrak q&\le3d-d_v+1.
 \end{split}
 \tag{G101.25}
\]
The argument uses the joint kernel \(K_w\) at its last step,
not the kernel of \(\gamma\).

\emph{7. The off-triangular fiber contributes at most one.}
Let \(\eta=\mathfrak p\cap A\).
Since \(\mathfrak p\) is off both triangular loci, some
positive-weight and some negative-weight entry survive in
the domain \(A/\eta\).
Put \(K=\operatorname{Frac}(B/\mathfrak q)\).
Choose surviving elements \(z,y\) of weights \(s>0,-t<0\).
Then \(z^t y^s\) is a nonzero element of \(K\).
Every homogeneous generator of weight \(b\) is algebraic
over \(K(z)\): for \(b<0\) use its \(s\)-th power times
\(z^{-b}\), which has weight zero; for \(b>0\) first use
its \(t\)-th power times \(y^b\), and then the preceding
relation for \(y\).  Weight-zero elements already belong
to \(B\).
Hence
\[
        \operatorname{trdeg}_{K}\operatorname{Frac}(A/\eta)
                                                        \le1.
 \tag{G101.26}
\]
This is the complete graded calculation of the literal source
\texttt{TrdegQuotLeOne.lean}, and it uses no single-character
residual identification.

We record the commutative-algebra interface for completion.
The ring \(B\) is complete Noetherian local over the complete
coefficient DVR and is a quotient of a power-series ring,
so it is excellent and universally catenary.
The dimension formula for a finite-type algebra at a closed
point over its maximal ideal, followed by the finite unramified
extension \(\mathcal O'/\mathcal O\), gives
\[
       \dim (A'_y/(\mathfrak p\cap A'_y))
          \le \dim B/\mathfrak q+
              \operatorname{trdeg}_{K}
                         \operatorname{Frac}(A/\eta).
 \tag{G101.27}
\]
Finite coefficient extension does not alter the latter
transcendence degree.
Completion preserves the dimension of this local quotient,
and \(C_x/\mathfrak p\) is a quotient of its completion.
Thus its dimension is bounded by the right side of
(G101.27).  These are exactly the generic
\texttt{dim\_gen} and \texttt{td\_fin} steps of
\texttt{LocDim.lean}; their hypotheses are supplied by
AC-G95 and (G101.4), not by an unproved old \texttt{Frame.Ok}
instantiation.
Combining (G101.25)--(G101.27) proves (G101.13).

Finally suppose \(T_+\not\subset x\).
The ideal \(T_+C_x\) contains a unit, so no prime of \(C_x\)
contains it.  A prime of the quotient in (G101.15) is either
off both triangular loci, where (G101.13) applies, or contains
\(T_-C_x\), where AC-G98 bounds its dimension by \(D_{\rm tri}\).
Taking the supremum over these primes proves (G101.15)
and its strict version.

\emph{8. Finite-image ratios, rather than ordinary adjacent equality.}
Let \(\mathfrak p\) contain \(\varpi\), and put
\(D=C_x/\mathfrak p\), a local domain of characteristic \(p\)
with coefficient field \(k'\).  A finite-order unit \(u\in D\)
with residue one must be one.  To see this, write its order
as \(p^a m\), with \(p\nmid m\).
The equality \((u^m-1)^{p^a}=0\) and the domain property
give \(u^m=1\); then
\(1+u+\cdots+u^{m-1}\) has the nonzero residue \(m\)
and is a unit, so \(u=1\).
Consequently reduction is injective on every finite subgroup
of \(D^\times\).

If \(\kappa_{v,i}/\kappa_{v,j}\) has finite image on \(N_v\),
each value is its fixed residual value in \(k^\times\);
divide by that coefficient-field value and use the preceding
argument.  This is precisely the vanishing of
\(\widehat{\mathfrak c}_{v,ij}\) in \(D\).
Conversely this vanishing makes the ratio its finite residual
character, so the image is finite.
For a continuous character on a profinite group, finite image
is equivalent to having an open kernel.  Thus outside this
locus no open subgroup of \(N_v\) can have trivial ratio.
Its intersection with an open subgroup of the full local group
is open in \(N_v\), proving the asserted stronger separation.
This checks the residual-cubic correction explicitly.
An adjacent ratio which is constantly its residual cubic
character has no ordinary equality collision but does belong
to this normalized bad locus.
No conclusion about the source's remaining global
\texttt{NiceOf} or propagation hypotheses is inferred here.
This completes the proof.
\end{proof}

\paragraph{Exact source interfaces and remaining scope.}
All local source files below are pinned at
\texttt{7cbf7a0c32fca1badf05931879a37c3f8a2b03b4},
under
\path{Development/AutomorphicArithmetic/AutomorphicVarieties/}.
The complete applied containing proofs read for this argument are
\path{FrobeniusSmoothVectors/DefIso.lean},
\path{LocDim.lean}, \path{TameInertiaCondition.lean},
\path{TrdegQuotLeOne.lean},
\path{KrullDimLtOfNotLeTriangular.lean},
\path{PresDefs.lean}, \path{PrGood.lean},
\path{PrCount.lean}, \path{PrBdim.lean}, \path{PrCfix.lean},
\path{PrDual.lean}, \path{PrBrig.lean}, \path{PrBstab.lean},
\path{PrSetlift.lean}, \path{BorelProjectionIsTwoCocycle.lean},
\path{PrCocv.lean}, \path{PrCorrect.lean},
\path{PrNorm1.lean}, \path{PrNorm.lean},
\path{PrTsurj.lean}, \path{HasObs.lean}, and
\path{PrMazur.lean}.
The literal declarations are \texttt{def\_iso},
\texttt{loc\_dim}, \texttt{l52}, \texttt{fibre},
\texttt{pr\_count}, \texttt{pr\_tsurj}, and
\texttt{Loc.Pres}; the sidecar manifest records each exact
URL, file hash and the actual reading scope.
The original statements retain their original residual hypothesis.
We use the newly proved completion lemma, the actual local dual
calculation, and the complete relation argument above to make
the changed-tuple transfer explicit.
The deeper representability and arithmetic duality inputs remain
those already identified in AC-G94 and AC-G96.
No Lean execution is claimed.

This proves the ordinary presentation and the local
finite-image-ratio dimension cut, with residual normalization.
It does not yet prove the source's stronger
\texttt{Loc.EPres}, \texttt{Sharp.chained}, or
\texttt{Sharp.dich} assertions for the changed tuple.
The actual adjacent-sector upper certificate in AC-G98,
and the final native support/curve propagation argument, remain
separate gates.  The improved local threshold \(d_v>2\) is
an exact saving in this collision cut; it is not an assertion
that every other old numerical hypothesis can be relaxed.

Geometrically, \(C_x\) is the formal completion of the actual
graded representation chart at its residual local system,
and (G101.23) is its ordinary polarized deformation obstruction
space.  The local character differences cut the repeated-character
subschemes of this formal chart.  This uses the established
finite-etale/local-system and complete-local deformation dictionary;
it attributes no new Grothendieck archive reading.
